\documentclass[12pt, letterpaper]{article}

% ── Paquetes ──────────────────────────────────────────────────────────────────
\usepackage[utf8]{inputenc}
\usepackage[T1]{fontenc}
\usepackage[spanish]{babel}
\usepackage{geometry}
\usepackage{graphicx}
\usepackage{hyperref}
\usepackage{listings}
\usepackage{xcolor}
\usepackage{booktabs}
\usepackage{tabularx}
\usepackage{array}
\usepackage{fancyhdr}
\usepackage{titlesec}
\usepackage{enumitem}
\usepackage{tikz}
\usepackage{float}
\usepackage{microtype}
\usepackage{parskip}
\usepackage{mdframed}
\usetikzlibrary{shapes.geometric, arrows.meta, positioning}

% ── Geometría de página ───────────────────────────────────────────────────────
\geometry{
    top=2.5cm,
    bottom=2.5cm,
    left=3cm,
    right=3cm
}

% ── Colores ───────────────────────────────────────────────────────────────────
\definecolor{cgsp-blue}{RGB}{15, 76, 129}
\definecolor{cgsp-gray}{RGB}{240, 242, 245}
\definecolor{cgsp-green}{RGB}{39, 174, 96}
\definecolor{cgsp-red}{RGB}{192, 57, 43}
\definecolor{cgsp-orange}{RGB}{211, 84, 0}
\definecolor{code-bg}{RGB}{248, 249, 250}
\definecolor{code-comment}{RGB}{106, 153, 85}
\definecolor{code-keyword}{RGB}{86, 156, 214}
\definecolor{code-string}{RGB}{206, 145, 120}

% ── Estilo de listings (bloques de código) ────────────────────────────────────
\lstdefinestyle{cgsp-proto}{
    backgroundcolor=\color{code-bg},
    basicstyle=\ttfamily\small,
    breaklines=true,
    breakatwhitespace=false,
    frame=single,
    framesep=4pt,
    rulecolor=\color{cgsp-blue!40},
    xleftmargin=8pt,
    xrightmargin=8pt,
    showstringspaces=false,
    tabsize=2,
    commentstyle=\color{code-comment},
    keywordstyle=\color{code-keyword}\bfseries,
    stringstyle=\color{code-string},
    numbers=none,
    captionpos=b,
    keepspaces=true
}

\lstset{style=cgsp-proto}

% ── Encabezado y pie de página ────────────────────────────────────────────────
\pagestyle{fancy}
\fancyhf{}
\fancyhead[L]{\small\textcolor{cgsp-blue}{\textbf{RFC-CGSP-1}}}
\fancyhead[C]{\small CyberGame Socket Protocol v1.0}
\fancyhead[R]{\small Abril 2026}
\fancyfoot[C]{\thepage}
\renewcommand{\headrulewidth}{0.5pt}
\renewcommand{\headrule}{\hbox to\headwidth{\color{cgsp-blue}\leaders\hrule height \headrulewidth\hfill}}

% ── Títulos de sección ────────────────────────────────────────────────────────
\titleformat{\section}
    {\large\bfseries\color{cgsp-blue}}
    {\thesection.}{0.5em}{}[\titlerule]

\titleformat{\subsection}
    {\normalsize\bfseries\color{cgsp-blue!80}}
    {\thesubsection.}{0.5em}{}

\titleformat{\subsubsection}
    {\normalsize\bfseries\color{cgsp-blue!60}}
    {\thesubsubsection.}{0.5em}{}

% ── Cuadro de nota ────────────────────────────────────────────────────────────
\newmdenv[
    backgroundcolor=cgsp-gray,
    linecolor=cgsp-blue,
    linewidth=2pt,
    leftline=true,
    rightline=false,
    topline=false,
    bottomline=false,
    innerleftmargin=10pt,
    innerrightmargin=10pt,
    innertopmargin=6pt,
    innerbottommargin=6pt
]{nota}

% ── Hipervínculos ─────────────────────────────────────────────────────────────
\hypersetup{
    colorlinks=true,
    linkcolor=cgsp-blue,
    urlcolor=cgsp-blue,
    citecolor=cgsp-blue
}

% ─────────────────────────────────────────────────────────────────────────────
%  DOCUMENTO
% ─────────────────────────────────────────────────────────────────────────────
\begin{document}

% ── Portada ───────────────────────────────────────────────────────────────────
\begin{titlepage}
    \centering
    \vspace*{1cm}

    {\small\texttt{Especificación de Protocolo}}\\[0.3cm]
    {\Huge\bfseries\color{cgsp-blue} RFC-CGSP-1}\\[0.5cm]
    {\large\bfseries CyberGame Socket Protocol (CGSP)}\\[0.2cm]
    {\large Versión 1.0}

    \vspace{1.5cm}
    \rule{\textwidth}{1pt}\\[0.4cm]
    {\normalsize
        \begin{tabular}{ll}
            \textbf{Estado:}   & Propuesto \\[4pt]
            \textbf{Fecha:}    & Abril 2026 \\[4pt]
            \textbf{Materia:}  & Internet: Arquitectura y Protocolos --- 2026-1 \\[4pt]
            \textbf{Autores:}  & Equipo Entrega\_1\_Telematica
        \end{tabular}
    }\\[0.4cm]
    \rule{\textwidth}{1pt}

    \vspace{1.5cm}
    \begin{mdframed}[backgroundcolor=cgsp-gray, linecolor=cgsp-blue, linewidth=1pt]
        \textbf{Abstract.} Este documento especifica el \textit{CyberGame Socket Protocol
        (CGSP)}, un protocolo de la capa de aplicación diseñado para soportar la
        comunicación en tiempo real entre clientes y el servidor central de un simulador
        interactivo multijugador de ciberseguridad. El protocolo define el formato de
        mensajes, el vocabulario de comandos, las reglas de procedimiento y el manejo de
        errores para una partida en la que jugadores con roles de ``Atacante'' o
        ``Defensor'' interactúan sobre un plano virtual compartido.
    \end{mdframed}

    \vfill
    {\small Universidad EAFIT --- Departamento de Ingeniería de Sistemas}
\end{titlepage}

% ── Tabla de contenidos ───────────────────────────────────────────────────────
\tableofcontents
\newpage

% ═════════════════════════════════════════════════════════════════════════════
\section{Visión General del Protocolo}
% ═════════════════════════════════════════════════════════════════════════════

\subsection{Propósito}

CGSP permite que múltiples jugadores interactúen simultáneamente con un servidor de
juego central. A través de este protocolo, los jugadores pueden:

\begin{itemize}[leftmargin=2em]
    \item Autenticarse en el sistema mediante un servicio de identidad externo.
    \item Listar y unirse a salas de juego activas.
    \item Moverse por un plano virtual bidimensional.
    \item Ejecutar acciones específicas de su rol (atacar o defender recursos críticos).
    \item Recibir notificaciones de eventos del sistema en tiempo real.
\end{itemize}

\subsection{Modelo de Funcionamiento}

CGSP adopta el modelo \textbf{cliente-servidor} con comunicación \textbf{bidireccional
persistente}. El siguiente diagrama ilustra la arquitectura del sistema:

\begin{figure}[H]
\centering
\begin{tikzpicture}[
    node distance=3.5cm,
    box/.style={rectangle, draw=cgsp-blue, fill=cgsp-gray, rounded corners=4pt,
                minimum width=3cm, minimum height=1cm, align=center, font=\small},
    arrow/.style={-{Stealth[length=6pt]}, thick, color=cgsp-blue}
]
    \node[box] (cliente) {Cliente\\(Jugador)};
    \node[box, right=of cliente] (servidor) {Servidor\\Central (C)};
    \node[box, right=of servidor] (identity) {Servicio de\\Identidad};

    \draw[arrow] ([yshift=3pt]cliente.east) -- ([yshift=3pt]servidor.west)
        node[midway, above, font=\tiny] {TCP persistente};
    \draw[arrow] ([yshift=-3pt]servidor.west) -- ([yshift=-3pt]cliente.east)
        node[midway, below, font=\tiny] {CGSP texto/\textbackslash{}n};

    \draw[arrow] ([yshift=3pt]servidor.east) -- ([yshift=3pt]identity.west)
        node[midway, above, font=\tiny] {TCP req/resp};
    \draw[arrow] ([yshift=-3pt]identity.west) -- ([yshift=-3pt]servidor.east)
        node[midway, below, font=\tiny] {OK ROLE / ERR};
\end{tikzpicture}
\caption{Arquitectura del sistema CGSP}
\end{figure}

El servidor acepta conexiones TCP entrantes y crea un hilo POSIX (\texttt{pthread})
dedicado por cada cliente. La conexión se mantiene activa durante toda la sesión del
jugador. El servidor también actúa como cliente frente al Servicio de Identidad para
autenticar usuarios.

\subsection{Capa de la Arquitectura TCP/IP}

CGSP opera en la \textbf{capa de aplicación} del modelo TCP/IP:

\begin{table}[H]
\centering
\begin{tabular}{>{\bfseries}l l}
\toprule
Capa & Protocolo \\
\midrule
Aplicación  & \textbf{CGSP v1.0} (este documento) \\
Transporte  & TCP (SOCK\_STREAM) \\
Red         & IPv4 \\
Enlace      & Ethernet / Wi-Fi \\
\bottomrule
\end{tabular}
\caption{Pila de protocolos de CGSP}
\end{table}

\begin{nota}
\textbf{Justificación de TCP sobre UDP.} El juego requiere confiabilidad absoluta en
la entrega de mensajes. Un paquete de \texttt{ATTACK} o \texttt{DEFEND} que se
pierda comprometería la consistencia del estado del juego entre servidor y clientes.
TCP garantiza entrega ordenada, detección de pérdidas y retransmisión automática.
UDP sería adecuado solo para actualizaciones de posición de alta frecuencia donde
una pérdida es aceptable --- lo cual no aplica a este diseño.
\end{nota}

\subsection{Resolución de Nombres}

CGSP prohíbe el uso de direcciones IP codificadas (\textit{hardcoded}) en el
software. Toda referencia a servicios externos (servidor de juego, servicio de
identidad) se resuelve mediante DNS usando la función \texttt{getaddrinfo()} de
POSIX. Si la resolución falla, el servicio maneja la excepción sin finalizar su
ejecución, respondiendo al cliente con \texttt{ERR 503}.

% ═════════════════════════════════════════════════════════════════════════════
\section{Especificación del Servicio}
% ═════════════════════════════════════════════════════════════════════════════

\subsection{Primitivas del Servicio}

Las primitivas de CGSP se organizan según el estado del cliente en el servidor.

\subsubsection{Estado: CONECTADO (antes de autenticarse)}

\begin{table}[H]
\centering
\begin{tabularx}{\textwidth}{l l X}
\toprule
\textbf{Primitiva} & \textbf{Dirección} & \textbf{Descripción} \\
\midrule
\texttt{AUTH} & Cliente $\rightarrow$ Servidor & Autenticar un usuario contra el servicio de identidad \\
\texttt{QUIT} & Cliente $\rightarrow$ Servidor & Cerrar la conexión ordenadamente \\
\bottomrule
\end{tabularx}
\caption{Primitivas disponibles en estado CONECTADO}
\end{table}

\subsubsection{Estado: AUTENTICADO}

\begin{table}[H]
\centering
\begin{tabularx}{\textwidth}{l l X}
\toprule
\textbf{Primitiva} & \textbf{Dirección} & \textbf{Descripción} \\
\midrule
\texttt{LIST\_ROOMS}  & Cliente $\rightarrow$ Servidor & Obtener lista de salas activas \\
\texttt{CREATE\_ROOM} & Cliente $\rightarrow$ Servidor & Crear una nueva sala de juego \\
\texttt{JOIN}         & Cliente $\rightarrow$ Servidor & Unirse a una sala existente \\
\texttt{STATUS}       & Cliente $\rightarrow$ Servidor & Consultar estado actual del jugador \\
\texttt{QUIT}         & Cliente $\rightarrow$ Servidor & Cerrar la conexión \\
\bottomrule
\end{tabularx}
\caption{Primitivas disponibles en estado AUTENTICADO}
\end{table}

\subsubsection{Estado: EN PARTIDA}

\begin{table}[H]
\centering
\begin{tabularx}{\textwidth}{l l l X}
\toprule
\textbf{Primitiva} & \textbf{Dirección} & \textbf{Rol} & \textbf{Descripción} \\
\midrule
\texttt{MOVE}   & C $\rightarrow$ S & Ambos    & Mover al jugador en el plano \\
\texttt{SCAN}   & C $\rightarrow$ S & ATTACKER & Detectar recursos críticos cercanos \\
\texttt{ATTACK} & C $\rightarrow$ S & ATTACKER & Lanzar un ataque sobre un recurso \\
\texttt{DEFEND} & C $\rightarrow$ S & DEFENDER & Mitigar un ataque activo \\
\texttt{START}  & C $\rightarrow$ S & Ambos    & Intentar iniciar la partida \\
\texttt{STATUS} & C $\rightarrow$ S & Ambos    & Consultar estado del jugador/sala \\
\texttt{QUIT}   & C $\rightarrow$ S & Ambos    & Abandonar la partida \\
\bottomrule
\end{tabularx}
\caption{Primitivas disponibles EN PARTIDA}
\end{table}

\subsubsection{Mensajes Asíncronos del Servidor}

\begin{table}[H]
\centering
\begin{tabularx}{\textwidth}{l l X}
\toprule
\textbf{Evento} & \textbf{Receptores} & \textbf{Descripción} \\
\midrule
\texttt{EVENT PLAYER\_JOINED}  & Sala completa    & Un nuevo jugador se unió a la sala \\
\texttt{EVENT PLAYER\_LEFT}    & Sala completa    & Un jugador abandonó la sala \\
\texttt{EVENT GAME\_STARTED}   & Sala completa    & La partida comenzó \\
\texttt{EVENT RESOURCE\_INFO}  & Solo defensores  & Coordenadas de recursos críticos (al iniciar) \\
\texttt{EVENT RESOURCE\_FOUND} & Solo atacante    & Recurso detectado por SCAN \\
\texttt{EVENT ATTACK}          & Sala completa    & Notificación de ataque activo \\
\texttt{EVENT DEFENDED}        & Sala completa    & Recurso fue defendido exitosamente \\
\texttt{EVENT GAME\_OVER}      & Sala completa    & Fin de partida con resultado \\
\bottomrule
\end{tabularx}
\caption{Mensajes asíncronos del servidor (broadcast y unicast)}
\end{table}

% ═════════════════════════════════════════════════════════════════════════════
\section{Formato de Mensajes}
% ═════════════════════════════════════════════════════════════════════════════

\subsection{Codificación General}

Todo mensaje CGSP debe cumplir las siguientes reglas de codificación:

\begin{itemize}[leftmargin=2em]
    \item \textbf{Codificación de caracteres:} UTF-8.
    \item \textbf{Delimitador de mensaje:} \texttt{\textbackslash{}n} (LF, \texttt{0x0A}). Un mensaje por línea.
    \item \textbf{Separador de campos:} espacio (\texttt{0x20}).
    \item \textbf{Longitud máxima de mensaje:} 2\,048 bytes.
    \item \textbf{Sensibilidad a mayúsculas en comandos:} NO (el servidor normaliza a mayúsculas).
\end{itemize}

\textbf{Estructura general:}
\begin{lstlisting}
COMANDO [PARAM1] [PARAM2] ... [PARAMn]\n
\end{lstlisting}

\subsection{Mensajes del Cliente hacia el Servidor}

\subsubsection{AUTH --- Autenticación}

\begin{lstlisting}
AUTH <usuario> <password>\n
\end{lstlisting}

\begin{table}[H]
\centering
\begin{tabularx}{\textwidth}{l l l X}
\toprule
\textbf{Campo} & \textbf{Tipo} & \textbf{Máx.} & \textbf{Descripción} \\
\midrule
\texttt{usuario}  & string & 63 chars & Nombre de usuario registrado \\
\texttt{password} & string & 63 chars & Contraseña del usuario \\
\bottomrule
\end{tabularx}
\end{table}

\textbf{Ejemplo:} \texttt{AUTH carlos hack2026}

\subsubsection{LIST\_ROOMS --- Listar Salas}

\begin{lstlisting}
LIST_ROOMS\n
\end{lstlisting}

Sin parámetros. Solicita la lista de salas de juego activas.

\subsubsection{CREATE\_ROOM --- Crear Sala}

\begin{lstlisting}
CREATE_ROOM\n
\end{lstlisting}

Sin parámetros. Crea una nueva sala de juego y coloca los recursos críticos en posiciones aleatorias del mapa.

\subsubsection{JOIN --- Unirse a Sala}

\begin{lstlisting}
JOIN <room_id>\n
\end{lstlisting}

\begin{table}[H]
\centering
\begin{tabularx}{\textwidth}{l l X}
\toprule
\textbf{Campo} & \textbf{Tipo} & \textbf{Descripción} \\
\midrule
\texttt{room\_id} & entero positivo & Identificador numérico de la sala \\
\bottomrule
\end{tabularx}
\end{table}

\textbf{Ejemplo:} \texttt{JOIN 42}

\subsubsection{START --- Iniciar Partida}

\begin{lstlisting}
START\n
\end{lstlisting}

Solicita el inicio de la partida. Solo tiene efecto si hay al menos un atacante y un defensor en la sala.

\subsubsection{MOVE --- Mover Jugador}

\begin{lstlisting}
MOVE <dx> <dy>\n
\end{lstlisting}

\begin{table}[H]
\centering
\begin{tabularx}{\textwidth}{l l l X}
\toprule
\textbf{Campo} & \textbf{Tipo} & \textbf{Rango} & \textbf{Descripción} \\
\midrule
\texttt{dx} & entero & {[}$-100$, $100${]} & Desplazamiento en el eje X \\
\texttt{dy} & entero & {[}$-100$, $100${]} & Desplazamiento en el eje Y \\
\bottomrule
\end{tabularx}
\end{table}

El servidor clampea la posición resultante dentro del mapa $[0, 99] \times [0, 99]$. \\
\textbf{Ejemplo:} \texttt{MOVE 5 -3}

\subsubsection{SCAN --- Escanear Área}

\begin{lstlisting}
SCAN\n
\end{lstlisting}

Solo para jugadores con rol \texttt{ATTACKER}. Detecta recursos críticos en un radio de \textbf{5 unidades} alrededor de la posición actual del atacante. La detección se calcula con distancia euclidiana:
\[
    d = \sqrt{(x_\text{jugador} - x_\text{recurso})^2 + (y_\text{jugador} - y_\text{recurso})^2} \leq 5
\]

\subsubsection{ATTACK --- Atacar Recurso}

\begin{lstlisting}
ATTACK <resource_id>\n
\end{lstlisting}

\begin{table}[H]
\centering
\begin{tabularx}{\textwidth}{l l X}
\toprule
\textbf{Campo} & \textbf{Tipo} & \textbf{Descripción} \\
\midrule
\texttt{resource\_id} & entero (0 ó 1) & Identificador del recurso crítico a atacar \\
\bottomrule
\end{tabularx}
\end{table}

\textbf{Ejemplo:} \texttt{ATTACK 0}

\subsubsection{DEFEND --- Defender Recurso}

\begin{lstlisting}
DEFEND <resource_id>\n
\end{lstlisting}

\begin{table}[H]
\centering
\begin{tabularx}{\textwidth}{l l X}
\toprule
\textbf{Campo} & \textbf{Tipo} & \textbf{Descripción} \\
\midrule
\texttt{resource\_id} & entero (0 ó 1) & Identificador del recurso a defender \\
\bottomrule
\end{tabularx}
\end{table}

\textbf{Ejemplo:} \texttt{DEFEND 0}

\subsubsection{STATUS --- Consultar Estado}

\begin{lstlisting}
STATUS\n
\end{lstlisting}

\subsubsection{QUIT --- Cerrar Sesión}

\begin{lstlisting}
QUIT\n
\end{lstlisting}

\subsection{Mensajes del Servidor hacia el Cliente}

\subsubsection{OK --- Operación Exitosa}

\begin{lstlisting}
OK <mensaje>\n
\end{lstlisting}

\textbf{Ejemplos:}
\begin{lstlisting}
OK Bienvenido carlos
OK Posicion: 15 42
OK Ataque lanzado
OK Recurso defendido
OK Hasta luego
\end{lstlisting}

\subsubsection{ROLE --- Rol Asignado}

\begin{lstlisting}
ROLE <ATTACKER|DEFENDER>\n
\end{lstlisting}

Enviado inmediatamente después de un \texttt{AUTH} exitoso.

\subsubsection{ERR --- Error}

\begin{lstlisting}
ERR <codigo> <descripcion>\n
\end{lstlisting}

\begin{table}[H]
\centering
\begin{tabularx}{\textwidth}{l l X}
\toprule
\textbf{Código} & \textbf{Nombre} & \textbf{Descripción} \\
\midrule
\texttt{400} & BAD\_REQUEST & Formato de mensaje inválido o parámetros faltantes \\
\texttt{401} & UNAUTHORIZED & Autenticación fallida o acción sin autenticar previamente \\
\texttt{403} & FORBIDDEN & El rol del jugador no permite esta acción \\
\texttt{404} & NOT\_FOUND & Sala o recurso no encontrado \\
\texttt{409} & CONFLICT & Conflicto de estado (sala llena, partida iniciada, recurso ya atacado) \\
\texttt{500} & INTERNAL\_ERROR & Error interno del servidor \\
\texttt{503} & SERVICE\_UNAVAILABLE & Servicio de identidad no disponible \\
\bottomrule
\end{tabularx}
\caption{Tabla de códigos de error de CGSP}
\end{table}

\textbf{Ejemplos:}
\begin{lstlisting}
ERR 401 Credenciales incorrectas
ERR 403 Solo los atacantes pueden usar ATTACK
ERR 404 Sala no encontrada
ERR 409 Recurso ya bajo ataque
ERR 503 Servicio de identidad no disponible, intenta luego
\end{lstlisting}

\subsubsection{ROOM\_CREATED --- Sala Creada}

\begin{lstlisting}
ROOM_CREATED <room_id>\n
\end{lstlisting}

\textbf{Ejemplo:} \texttt{ROOM\_CREATED 7}

\subsubsection{ROOM\_LIST --- Lista de Salas}

\begin{lstlisting}
ROOM_LIST <n>\n
ROOM <id> <WAITING|RUNNING|FINISHED> <jugadores>/<max>\n
...
\end{lstlisting}

\textbf{Ejemplo:}
\begin{lstlisting}
ROOM_LIST 2
ROOM 1 WAITING 1/10
ROOM 5 RUNNING 3/10
\end{lstlisting}

\subsubsection{Eventos del Servidor}

\begin{lstlisting}
EVENT PLAYER_JOINED <usuario> <ATTACKER|DEFENDER>\n
EVENT PLAYER_LEFT <usuario>\n
EVENT GAME_STARTED\n
EVENT RESOURCE_INFO <id> <x> <y>\n
EVENT RESOURCE_FOUND <id> <x> <y>\n
EVENT ATTACK <resource_id> <atacante>\n
EVENT DEFENDED <resource_id> <defensor>\n
EVENT GAME_OVER <ATTACKER|DEFENDER>\n
\end{lstlisting}

\begin{nota}
\textbf{Unicast vs.\ Broadcast.} Los eventos \texttt{RESOURCE\_INFO} se envían
únicamente al jugador defensor que se acaba de conectar. Los eventos \texttt{RESOURCE\_FOUND}
se envían únicamente al atacante que ejecutó el \texttt{SCAN}. Todos los demás
eventos se difunden a todos los jugadores de la sala (\textit{broadcast}).
\end{nota}

% ═════════════════════════════════════════════════════════════════════════════
\section{Reglas de Procedimiento (Flujo de Operación)}
% ═════════════════════════════════════════════════════════════════════════════

\subsection{Máquina de Estados del Cliente}

\begin{figure}[H]
\centering
\begin{tikzpicture}[
    node distance=2.2cm,
    state/.style={rectangle, draw=cgsp-blue, fill=cgsp-gray, rounded corners=6pt,
                  minimum width=5cm, minimum height=1cm, align=center, font=\small\bfseries},
    arrow/.style={-{Stealth[length=6pt]}, thick, color=cgsp-blue!70},
    label/.style={font=\tiny, color=cgsp-blue!80, align=center}
]
    \node[state] (S1) {CONECTADO\\{\footnotesize AUTH, QUIT}};
    \node[state, below=of S1] (S2) {AUTENTICADO\\{\footnotesize LIST\_ROOMS, CREATE\_ROOM, JOIN, STATUS, QUIT}};
    \node[state, below=of S2] (S3) {EN SALA\\{\footnotesize START, STATUS, QUIT}};
    \node[state, below=of S3] (S4) {EN PARTIDA\\{\footnotesize MOVE, SCAN/ATTACK, DEFEND, STATUS, QUIT}};

    \draw[arrow] (S1) -- (S2)
        node[midway, right, label] {AUTH exitoso\\$\rightarrow$ OK + ROLE};
    \draw[arrow] (S2) -- (S3)
        node[midway, right, label] {JOIN exitoso\\$\rightarrow$ OK};
    \draw[arrow] (S3) -- (S4)
        node[midway, right, label] {START ($\geq$1 ATC + $\geq$1 DEF)\\$\rightarrow$ GAME\_STARTED + RESOURCE\_INFO};
\end{tikzpicture}
\caption{Máquina de estados del cliente CGSP}
\end{figure}

\subsection{Ciclo de Vida de una Sala}

\begin{figure}[H]
\centering
\begin{tikzpicture}[
    node distance=3.5cm,
    state/.style={rectangle, draw=cgsp-blue, fill=cgsp-gray, rounded corners=4pt,
                  minimum width=2.8cm, minimum height=0.9cm, align=center, font=\small\bfseries},
    arrow/.style={-{Stealth[length=6pt]}, thick, color=cgsp-blue!70}
]
    \node[state] (W) {WAITING};
    \node[state, right=of W] (R) {RUNNING};
    \node[state, right=of R] (F) {FINISHED};

    \draw[arrow] (W) -- (R) node[midway, above, font=\tiny] {START exitoso};
    \draw[arrow] (R) -- (F) node[midway, above, font=\tiny] {condición\\de victoria};
\end{tikzpicture}
\caption{Ciclo de vida de una sala de juego}
\end{figure}

\subsection{Flujo de Ataque y Defensa}

\begin{lstlisting}
ATACANTE          SERVIDOR             DEFENSOR(ES)
   |                  |                     |
   |-- ATTACK 0 ----->|                     |
   |<- OK Ataque -----|                     |
   |                  |-- EVENT ATTACK 0 -->|
   |                  |                     |-- DEFEND 0 -->|
   |                  |<-------------------- DEFEND 0 ------|
   |                  |-- EVENT DEFENDED 0 maria ---------->|
   |<- EVENT DEFENDED |                     |
\end{lstlisting}

\begin{nota}
\textbf{Timer de ataque.} Un recurso bajo ataque que no sea defendido dentro de
\textbf{30 segundos} constituye victoria para el equipo atacante.
\end{nota}

\subsection{Flujo de Exploración del Atacante}

\begin{lstlisting}
ATACANTE              SERVIDOR
   |                     |
   |-- MOVE 30 35 ------>|   pos = (30, 35)
   |<- OK Posicion 30 35 |
   |-- SCAN ------------>|   dist((30,35), recurso) <= 5 ?
   |<- EVENT RESOURCE_FOUND 0 32 37    (si detecta)
   |<- OK SCAN: no encontrado          (si no detecta)
   |-- ATTACK 0 -------->|
   |<- OK Ataque lanzado |
\end{lstlisting}

\subsection{Manejo de Errores y Excepciones}

\begin{table}[H]
\centering
\begin{tabularx}{\textwidth}{X X}
\toprule
\textbf{Situación} & \textbf{Respuesta del Servidor} \\
\midrule
Línea vacía o sin \texttt{\textbackslash{}n} & Ignorada silenciosamente \\
Comando desconocido & \texttt{ERR 400 Comando desconocido} \\
Parámetros faltantes & \texttt{ERR 400 Uso: CMD <params>} \\
Acción sin autenticar & \texttt{ERR 401 Debes autenticarte primero} \\
Rol incorrecto para acción & \texttt{ERR 403 Solo los [rol] pueden usar [CMD]} \\
Sala inexistente & \texttt{ERR 404 Sala no encontrada} \\
Recurso inexistente & \texttt{ERR 404 Recurso no encontrado} \\
Sala llena o partida iniciada & \texttt{ERR 409 <descripción del conflicto>} \\
Servicio de identidad caído & \texttt{ERR 503} --- servidor \textbf{no} finaliza \\
Desconexión abrupta del cliente & \texttt{read()} retorna $\leq 0$ $\rightarrow$ cierre limpio del hilo \\
AUTH con sesión ya activa & \texttt{ERR 403 Ya estas autenticado} \\
JOIN cuando ya está en sala & \texttt{ERR 403 Ya estas en una sala} \\
\bottomrule
\end{tabularx}
\caption{Tabla de manejo de excepciones del protocolo}
\end{table}

% ═════════════════════════════════════════════════════════════════════════════
\section{Ejemplos de Implementación}
% ═════════════════════════════════════════════════════════════════════════════

\subsection{Sesión Completa de un Atacante}

\begin{lstlisting}
# Conexion establecida
S->C: OK Bienvenido al servidor CyberGame (CGSP v1.0). Usa AUTH

# Autenticacion
C->S: AUTH carlos hack2026
S->C: OK Bienvenido carlos
S->C: ROLE ATTACKER

# Listar salas disponibles
C->S: LIST_ROOMS
S->C: ROOM_LIST 1
S->C: ROOM 3 WAITING 1/10

# Unirse a sala
C->S: JOIN 3
S->C: OK Te uniste a la sala

# Iniciar partida
C->S: START
S->C: OK Partida iniciada
S->C: EVENT GAME_STARTED

# Explorar el mapa
C->S: MOVE 30 35
S->C: OK Posicion: 30 35

C->S: SCAN
S->C: EVENT RESOURCE_FOUND 0 32 37

# Lanzar ataque
C->S: ATTACK 0
S->C: OK Ataque lanzado

# Notificacion de defensa exitosa
S->C: EVENT DEFENDED 0 maria

C->S: QUIT
S->C: OK Hasta luego
\end{lstlisting}

\subsection{Sesión Completa de un Defensor}

\begin{lstlisting}
C->S: AUTH maria seg2026
S->C: OK Bienvenido maria
S->C: ROLE DEFENDER

C->S: JOIN 3
S->C: OK Te uniste a la sala

C->S: START
S->C: EVENT GAME_STARTED
S->C: EVENT RESOURCE_INFO 0 32 37    # conoce los recursos
S->C: EVENT RESOURCE_INFO 1 75 82

C->S: MOVE 32 37
S->C: OK Posicion: 32 37

# Alerta de ataque
S->C: EVENT ATTACK 0 carlos

C->S: DEFEND 0
S->C: OK Recurso defendido

C->S: QUIT
S->C: OK Hasta luego
\end{lstlisting}

\subsection{Diagrama de Secuencia de una Partida Completa}

\begin{lstlisting}
Atacante        Servidor         Defensor
   |               |                |
   |-- AUTH ------>|<----- AUTH ----|
   |<- OK/ROLE ----|---- OK/ROLE -->|
   |               |                |
   |-- CREATE_ROOM->|               |
   |<- ROOM_CREATED|                |
   |               |                |
   |-- JOIN 5 ---->|<----- JOIN 5 --|
   |<- OK ----------|----- OK ----->|
   |               |                |
   |-- START ------>|<----- START --|
   |<- GAME_STARTED |-- GAME_STARTED>|
   |                |-- RESOURCE_INFO>|
   |               |                |
   |-- MOVE/SCAN -->|               |
   |<- FOUND -------|               |
   |               |                |
   |-- ATTACK 0 -->|-- EVENT ATTACK>|
   |<- OK ---------|<--- DEFEND 0 --|
   |<- DEFENDED ----|---- OK ------>|
\end{lstlisting}

\subsection{Prueba Manual con Telnet}

Para verificar el protocolo sin necesidad de un cliente completo:

\begin{lstlisting}[language=bash]
# Terminal 1: Iniciar el servicio de identidad stub
python3 server/identity_stub.py 9090

# Terminal 2: Iniciar el servidor de juego
export IDENTITY_HOST=localhost
export IDENTITY_PORT=9090
cd server && ./server 8080 server.log

# Terminal 3: Conectar como atacante
telnet localhost 8080

# Terminal 4: Conectar como defensor
telnet localhost 8080
\end{lstlisting}

% ═════════════════════════════════════════════════════════════════════════════
\section{Parámetros del Servidor}
% ═════════════════════════════════════════════════════════════════════════════

\begin{table}[H]
\centering
\begin{tabularx}{\textwidth}{l l l X}
\toprule
\textbf{Parámetro} & \textbf{Fuente} & \textbf{Default} & \textbf{Descripción} \\
\midrule
Puerto de escucha    & CLI: \texttt{./server <puerto>}   & ---         & Puerto TCP del servidor \\
Archivo de logs      & CLI: \texttt{./server <p> <log>}  & ---         & Ruta del archivo de logging \\
Host de identidad    & Env: \texttt{IDENTITY\_HOST}      & localhost   & Hostname del servicio de identidad \\
Puerto de identidad  & Env: \texttt{IDENTITY\_PORT}      & 9090        & Puerto TCP del servicio de identidad \\
\bottomrule
\end{tabularx}
\caption{Parámetros de configuración del servidor}
\end{table}

% ═════════════════════════════════════════════════════════════════════════════
\section{Parámetros del Mapa de Juego}
% ═════════════════════════════════════════════════════════════════════════════

\begin{table}[H]
\centering
\begin{tabular}{l r}
\toprule
\textbf{Parámetro} & \textbf{Valor} \\
\midrule
Dimensiones del mapa             & $100 \times 100$ unidades \\
Recursos críticos por sala       & 2 (posición aleatoria al crear sala) \\
Radio de detección (\textsc{SCAN}) & 5 unidades (distancia euclidiana) \\
Tiempo límite de defensa         & 30 segundos \\
Máx.\ jugadores por sala         & 10 \\
Máx.\ salas simultáneas          & 20 \\
\bottomrule
\end{tabular}
\caption{Parámetros del mapa del juego}
\end{table}

% ═════════════════════════════════════════════════════════════════════════════
\section{Consideraciones de Seguridad}
% ═════════════════════════════════════════════════════════════════════════════

\begin{itemize}[leftmargin=2em]
    \item Las contraseñas \textbf{no se almacenan} en el servidor de juego; se reenvían al
          servicio de identidad en cada autenticación.
    \item Las contraseñas \textbf{no se registran} en los logs del servidor (aparecen como
          \texttt{[***]}).
    \item No se implementa cifrado en esta versión (fuera del alcance del proyecto).
    \item Cada cliente opera en un hilo independiente; un cliente malicioso no puede
          bloquear a los demás.
\end{itemize}

% ─────────────────────────────────────────────────────────────────────────────
\vspace{2cm}
\hrule
\vspace{0.3cm}
{\small\centering\textit{Fin del documento RFC-CGSP-1 --- CyberGame Socket Protocol v1.0}\par}

\end{document}
